\chapter{Our workloads, concretely}
\label{sec:workloads}

The preceding chapters organized the literature by method. This one flips
the table and asks, for each family of work in our repositories, what the
evidence says we should actually do. Each section ends in a concrete
default: the study we would configure today, with the chapters that
justify each choice.

\section{Reinforcement learning}

The controlled evidence base here is unusually good, and we have already
leaned on it: \citet{eimer2023hyperparameters} tuned PPO, SAC, and DQN
across eleven environments with matched budgets and seeds. Four of their
findings should directly shape our defaults. First, hyperparameter
importance is broad but concentrated: almost every swept knob mattered
somewhere, importance analysis attributed most variance to one or two knobs
per environment, and which knobs those were changed across environments, so
the search space should include everything plausible rather than a
hand-picked pair. This is the effective-dimensionality picture of
Figure~\ref{fig:gridrandom} with a twist: the important axes exist, but we
cannot know in advance which they are, which is precisely the situation
where random-style coverage of every axis, and GP lengthscales that learn
to switch axes off (Section~\ref{sec:gp}), earn their keep.

Second, the
response surfaces themselves looked smooth and well behaved; the
difficulty of tuning RL comes overwhelmingly from seed noise, not from a
pathological objective. That is good news for every surrogate in
Chapter~\ref{sec:bayesian}, and it aims the engineering effort at noise
handling (the noise-aware acquisitions of
Section~\ref{sec:acquisition}) rather than at exotic models.

Third, the
seed noise is severe enough to corrupt conclusions: configurations tuned
on one seed transferred poorly, incumbents evaluated on tuning seeds
overstated test performance several-fold in the worst cases, and their
protocol of tuning on a handful of seeds and reporting on disjoint test
seeds is the minimum honest practice.

Fourth, at budgets of ten to
sixty-four full training runs, DEHB was the most dependable tuner, with
random search respectable and population methods behind; the population
methods' natural habitat is the large-parallelism regime of
Chapter~\ref{sec:population}, where BG-PBT's own experiments live.

Assembled into a default: an RL study below a few dozen full-run
equivalents uses DEHB (or BO plus ASHA) over the full plausible space,
evaluates each candidate on several tuning seeds, reports on disjoint
test seeds, and holds the tuning budget fixed in any method comparison;
a study with eight or more dedicated full-length workers and a
schedule-shaped question upgrades to the population line, PB2-Mix now
and BG-PBT when we build it.

The broader AutoRL literature, surveyed by
\citet{parkerholder2022autorl}, adds the framing that hyperparameters,
architecture, and even algorithm components form one design space, but the
budget-dependent boundary above is the piece with direct operational
consequence for us.

Two later results sharpen the picture. The ARLBench
benchmark, built from hyperparameter-landscape data across many
environments, reports that RL tuning landscapes are less benign than the
supervised-learning ones and selects small representative environment
subsets that cut the cost of evaluating a tuner by roughly an order of
magnitude, which is precisely what our validation suite needs
\citep{becktepe2024arlbench}. And a methodology paper on hyperparameter
sensitivity warns that some reported algorithmic improvements in RL amount
to increased reliance on tuning, one more reason to hold tuning budgets
fixed when comparing methods \citep{adkins2024sensitivity}.

For
meta-gradient methods that adapt hyperparameters inside a single run, such
as STAC \citep{zahavy2020selftuning}, our position is to treat them as
properties of the RL algorithm, not of the tuner: worth enabling when
RLlib exposes them, outside our product's scope to implement.

One more RL-specific decision deserves sunlight: what number the tuner
optimizes. Final evaluated return conflates stability with peak
performance; area under the training curve rewards fast learners; return
at a fixed step budget penalizes slow starters that would win given more
steps. The three orderings genuinely disagree, and
Figure~\ref{fig:curves} shows the disagreement in miniature: at the
final step the blue configuration wins, on area under the curve the red
one may, and a truncated budget flips the verdict again. The choice
changes rankings, it interacts with early stopping (a curve that
plateaus late looks like a curve worth killing), and it should be an
explicit, recorded field of every study rather than a hidden default.

When return and stability both matter, nothing stops us from declaring
them as two objectives and using the machinery of Chapter~\ref{sec:moo};
the multi-objective RL literature makes the matching argument on the
agent's side, warning against collapsing genuinely conflicting
objectives into a fixed linear combination \citep{hayes2022morl}, which
is Figure~\ref{fig:chebyshev}'s lesson wearing a different hat. We found
no dedicated literature on multi-objective hyperparameter tuning for RL
specifically; the generic machinery is the state of the art.

\section{Neural-transport samplers}

Our NeuTra-style work trains a normalizing flow to reshape a posterior,
then runs Hamiltonian Monte Carlo in the warped space
\citep{hoffman2019neutra}. Tuning has two loops, and the division of
labor matters enough to draw (Figure~\ref{fig:sampler-loops}).

The inner
loop belongs to the samplers themselves: step size adapts by the
dual-averaging scheme introduced with the No-U-Turn Sampler (NUTS),
trajectory length by NUTS's
recursive doubling \citep{hoffman2014nuts} or, in the
many-parallel-chains regime that accelerators favor, by gradient-based
adaptation of the trajectory length against a mixing proxy, as in
ChEES-HMC \citep{hoffman2021chees} and its successor that maximizes a
bound on effective sample size per gradient along an estimated principal
component \citep{sountsov2021snaper}.

Where good internal adaptation
exists, the outer tuner should not fight it: a tuner that overrides the
step size NUTS would have chosen is spending trials to re-derive
dual averaging, badly. What remains for the outer loop is everything the
samplers cannot adapt on their own: the flow's architecture (depth,
width, number of transformations), its optimizer and training length,
the variational objective's minibatch size, and the target acceptance
rate handed to the inner adapters.

\begin{figure}[t]
\centering
\begin{tikzpicture}[
  box/.style={draw, rounded corners, align=center, inner sep=5pt, font=\small},
  arr/.style={-{Latex}, thick},
  node distance=8mm and 24mm]
\node[box, minimum width=40mm, minimum height=15mm] (outer)
  {\textbf{outer tuner (ours)}\\flow depth, width, layers;\\optimizer, training length;\\target acceptance rate};
\node[box, right=of outer, minimum width=52mm, minimum height=15mm] (inner)
  {\textbf{inner adaptation (the sampler's)}\\step size: dual averaging\\trajectory length: NUTS / ChEES};
\draw[arr] (outer.east) -- node[above, font=\footnotesize] {configuration} (inner.west);
\draw[arr] (inner.east) -- ++(0.9,0) node[right, box, minimum width=30mm, align=center] (score)
  {ESS per gradient,\\after vetoes};
\draw[arr] (score.south) |- ++(-2.0,-0.75) -| node[below, pos=0.25, font=\footnotesize]
  {vetoes: divergences, split-$\widehat{R}$, reference checks} (outer.south);
\end{tikzpicture}
\caption{The two tuning loops of a neural-transport sampler. The inner
adapters already solve step size and trajectory length; the outer tuner
owns what they cannot reach, and its objective is efficiency (effective
samples per gradient) gated by correctness vetoes that no efficiency
score can buy off.}
\label{fig:sampler-loops}
\end{figure}

The outer objective already has a citable template, and it is the NeuTra
paper's own. Its authors selected step size and leapfrog steps by
Bayesian optimization on an objective that rewards effective sample size
(ESS, the number of independent draws a correlated chain is worth)
per gradient evaluation while sharply penalizing configurations whose
potential-scale-reduction diagnostic $\widehat{R}$ departs from one
\citep{hoffman2019neutra}; tuning samplers by BO goes back at least to
\citet{mahendran2012adaptive}. Effective samples per gradient is the
right currency because gradient evaluations dominate cost in both flow
training and leapfrog integration.

But ESS estimated from short chains
is exactly the kind of proxy that Chapter~\ref{sec:multifidelity} warned
about, and sampler tuning has a sharper failure mode than slow starters:
a chain that mixes quickly inside one mode of a multimodal posterior
reports excellent short-run ESS while being wrong, the statistical
equivalent of Figure~\ref{fig:curves}'s sprinter, except that this
sprinter never stalls visibly; it simply reports the wrong answer with
growing confidence.

The remedy is to treat correctness diagnostics as
vetoes rather than objectives: divergence counts, split $\widehat{R}$
across independent chains, and, where a reference posterior is available
on a pilot problem, agreement with it. A configuration that fails a veto
is discarded no matter how efficient it looks; among configurations that
pass, ESS per gradient is a legitimate objective. Multi-fidelity
structure is still available (shorter chains, fewer chains), but rung
promotion should require the vetoes to pass at every rung, which is
exactly the kind of user-defined gate our scheduler interfaces must
support (Chapter~\ref{sec:architecture}).

The concrete default: a single-objective study on ESS per gradient with
divergence and split-$\widehat{R}$ vetoes at every rung, ASHA over chain
length as the scheduler, the GP searcher of Chapter~\ref{sec:bayesian}
over the flow and training knobs, and the inner adapters left switched
on with only their target acceptance rate exposed to the search.

\section{Hamiltonian and physics-structured networks}

Hamiltonian neural networks learn a scalar function whose symplectic
gradient supplies the dynamics; conservation comes from the architecture
rather than from a penalty term, and the original paper reports models
that conserve an energy-like quantity where an unconstrained baseline
drifts \citep{greydanus2019hnn}. Notably for us, its authors chose their
hyperparameters by a coarse grid over learning rate, layer width,
activation, and batch size, which is precisely the procedure
Chapter~\ref{sec:modelfree} retired; the field's own tuning practice
trails its modeling practice, and that gap is our opportunity.

Tuning these models well is multi-objective in the evaluation even when
the training loss is a single term: short-horizon prediction error,
long-rollout state error, and energy drift move separately, the training
loss sees only the first, and the scientific value lives in the second
and third.

The axes of Figure~\ref{fig:pareto} were chosen with this
workload in mind, and the study template is the one that figure
illustrates: a two-or-three-objective search (prediction error,
long-rollout error or energy drift, optionally training cost) solved
with qLogNEHVI or MO-ASHA from Chapter~\ref{sec:moo}, over a space that
includes integrator choice and step size, width and depth, activation
(smoothness matters when the learned function is differentiated),
learning rate, and loss weights where present. Where variants do carry
multi-term losses (the pixel-observation task in the original paper
combines the dynamics loss with autoencoder terms), the weights joining
them are exactly the knobs an outer tuner should own.

Trials are cheap,
minutes rather than hours, so this is also the workload where the
tuner's own overhead (Section~\ref{sec:costmodel}) is most visible and
the local execution backend of Chapter~\ref{sec:architecture} earns its
keep.

A neighboring literature balances competing loss terms inside training
rather than outside it: gradient-statistics reweighting for
physics-informed networks \citep{wang2020pathologies}, weights set
from neural-tangent-kernel eigenvalues \citep{wang2020ntk}, and relative
loss balancing schemes compared systematically by
\citet{bischof2025relobralo}. The division-of-labor rule from the
sampler subsection transfers unchanged: where internal balancing works,
let the trainer balance and let the tuner search what the trainer
cannot, and when in doubt, the internal balancer is a competitor our
outer loop must beat at matched cost, not a rival to ignore.

\section{Architecture search, in moderation}

Our NAS folder is full of ambitious methods, and the sober reading of the
field is that a small team should adopt its cheap findings and skip its
expensive machinery. The thousand-paper survey of \citet{white2023nas1000}
recommends random search as the mandatory baseline for any NAS claim
(Chapter~\ref{sec:modelfree}'s discipline, enforced field-wide) and
records that simple quantities such as parameter count and FLOPs
(floating-point operation counts) are
surprisingly competitive with the leading zero-cost proxies, whose
reliability degrades on larger search spaces. Systematic bench-marking of
those proxies found their value to lie in complementing a learned
performance predictor rather than in standalone selection
\citep{krishnakumar2022nasbenchsuite}.

On the one-shot side, the DARTS
family is documented to produce degenerate architectures across a range
of search spaces without careful regularization \citep{zela2020robustifying},
and controlled evaluations found much of the reported progress
attributable to training tricks rather than to the search itself, with
hand-designed macro-structure mattering more than the searched cells
\citep{yang2020frustratingly}.

Meanwhile the benchmark built for joint
architecture-and-hyperparameter search motivates itself with the
observation that good hyperparameters can matter more than the best
architecture \citep{bansal2022jahs}.

For us, the conclusion is comfortable: architecture search should mean
width, depth, and a few structural flags included as ordinal and
categorical axes in the ordinary search space (as BG-PBT does for policy
and value networks, with the ordinal-aware kernels of
Section~\ref{sec:highdim} doing the modeling), multi-fidelity doing the
heavy lifting, parameter count reported as an objective when size
matters, and no dedicated weight-sharing subsystem. That covers every
architecture question our current projects actually ask, at a small
fraction of the complexity.

Because ``in moderation'' is the kind of phrase that drifts, the
boundary is worth stating as a scope rather than as a mood. Inside it:
width and depth as ordinal axes, a small number of structural flags
(normalization choice, residual connections, activation family) as
categorical axes, parameter count and measured FLOPs available as
objectives for the multi-objective machinery of
Chapter~\ref{sec:moo} rather than as proxies substituting for
training, and architecture coordinates mutating through the same
explore step every other coordinate uses, which is what BG-PBT's
architecture handling amounts to once the weight-sharing is stripped
out (Chapter~\ref{sec:population}). Outside it, deliberately: topology
and cell search, one-shot or weight-sharing supernets, a zero-cost
proxy subsystem, grammar-based or hierarchical architecture spaces
\citep{schrodi2023grammar}, and differentiable search of any flavor.
The exclusions are not a claim that those methods fail; they are the
claim this chapter's evidence actually supports, which is that they
cost a subsystem each and buy nothing our workloads have asked for. A
project that needs them needs a different governing document, and the
honest move then is to write one rather than to widen this sentence.

One consequence is worth flagging where it lands. Changing width or
depth mid-run invalidates the weights a population method wants to
inherit, so a population arm that mutates an architecture coordinate
pays for the change: either it restarts the trial from initialization,
or it transfers what it can by distillation from the parent, which is
the route BG-PBT takes \citep{wan2022bgpbt} and which
Chapter~\ref{sec:population} treats as part of the population line's
price rather than as free.

\section{The compute bill, and the shortcuts we are asked about}
\label{sec:computebill}

Tuning reinforcement learning is expensive enough that the fair question
is not ``which method'' but ``can we afford any of this,'' so the bill
deserves to be written out. A study's cost factors cleanly into three
multiplied terms: the number of candidate configurations examined, times
the seeds per candidate that Chapter~\ref{sec:workloads}'s protocol
demands, times the training steps per trial. In the running example the
last factor alone is up to 150 million environment steps, hours on
several accelerators, so a naive 64-candidate, 3-seed, full-length study
is roughly two hundred full training runs before anything has been
learned about tuning at all.

Every lever this document recommends
attacks one factor: multi-fidelity scheduling cuts steps per trial by
an order of magnitude (Chapter~\ref{sec:multifidelity}, and it is the
biggest single lever); model-based search and warm starts cut the
number of candidates (Chapters~\ref{sec:bayesian}
and~\ref{sec:transfer}); the seed protocol stops compute being spent
optimizing luck; population methods make the trials \emph{be} the
training run rather than a rehearsal for it
(Chapter~\ref{sec:population}); and ARLBench-style representative
benchmark subsets cut the cost of evaluating tuner changes by roughly
an order of magnitude \citep{becktepe2024arlbench}, which is what makes
our own validation suite affordable.

Against that backdrop, three architecture-search shortcuts come up
whenever the bill is discussed, and each deserves a plain description
and a verdict.

\emph{One-shot search} trains a single ``supernet'' that contains every
candidate architecture as a sub-network with shared weights; candidates
are then scored by inheriting weights from the supernet instead of being
trained separately, so one big training run amortizes thousands of
evaluations. The economics are attractive and the evidence is not. The
DARTS family of one-shot methods is documented to produce degenerate
architectures across search spaces without careful regularization
\citep{zela2020robustifying}; controlled re-evaluations attribute much
of the reported progress to training tricks rather than the search, with
hand-designed macro-structure mattering more than the searched cells
\citep{yang2020frustratingly}; and the field's own thousand-paper survey
responds by making random search the mandatory baseline
\citep{white2023nas1000}.

Two further points are specific to us rather
than to the literature, and we label them as our inference: supernet
training assumes a stable supervised training signal, while RL's data
distribution shifts under the policy as it learns; and our architecture
questions concern small policy and value networks, where there is not
enough training to amortize for the supernet's economy to pay. Verdict:
no one-shot subsystem, for reliability reasons first and fit reasons
second.

\emph{Zero-shot search} (zero-cost proxies) scores architectures at
initialization, without any training, from quantities like gradient
statistics at the first batch; the promise is architecture selection
for the price of a forward-backward pass. The systematic benchmark of
these proxies found their value to lie in complementing a learned
performance predictor rather than in standalone selection, with
reliability degrading on larger spaces
\citep{krishnakumar2022nasbenchsuite}, and the survey notes that plain
parameter count and FLOPs are surprisingly competitive with the
sophisticated proxies \citep{white2023nas1000}. That last fact is the
actionable one: params and FLOPs are free, we already record them in
the run store, and they can serve as objectives or tie-breakers today.
Verdict: report the free quantities; build no proxy subsystem.

\emph{Grammar-based spaces} attack a different problem than cost.
A context-free grammar generates architectures the way a grammar
generates sentences: production rules expand high-level structure into
blocks and blocks into operations, which compactly defines hierarchical
search spaces vastly larger than the usual cell catalogs, searched with
Bayesian optimization over hierarchical kernels
\citep{schrodi2023grammar}; NePS ships support for such spaces
(Chapter~\ref{sec:libraries}). The honest reading for a manager: this
is a solution to \emph{expressiveness}, letting one search whole
macro-structures rather than a fixed menu, and a larger space raises
rather than lowers the number of evaluations needed. It becomes
relevant the day a project genuinely needs to discover topology; every
architecture question our projects ask today is width, depth, and
flags, where the benchmark evidence says hyperparameters matter at
least as much as the architecture anyway \citep{bansal2022jahs}.
Verdict: not in the package now; NePS is the reference implementation
if a topology-search project ever lands.

The summary sentence for the section: the compute problem in RL tuning
is real, and its solution in our package is the stack of levers in the
first paragraph, not an architecture-search shortcut; the shortcuts
either lack reliability evidence (one-shot), reduce to free quantities
we already record (zero-shot), or solve a problem we do not have
(grammar spaces).

\begin{decisions}
\noindent\textbf{Decisions from this chapter.} Per workload, the study
we configure today; and the architecture-search verdicts.
\begin{itemize}[leftmargin=1.3em, itemsep=2pt]
\item \textbf{RL, routine budgets} (a few dozen full runs): DEHB or
BO-plus-ASHA over the full plausible space, several tuning seeds,
disjoint test seeds, objective recorded explicitly. \textbf{RL, large
parallelism} (8+ workers): the population line, PB2-Mix now, BG-PBT
when built.
\item \textbf{Neural-transport samplers}: GP searcher over flow and
training knobs, ASHA over chain length, ESS-per-gradient objective,
divergence and split-$\widehat{R}$ vetoes at every rung; leave the
samplers' internal adapters on.
\item \textbf{Hamiltonian networks}: two-to-three-objective study
(prediction error, energy drift or rollout error, optionally cost)
with qLogNEHVI or MO-ASHA, on the local execution backend.
\item \textbf{Architecture search}: width, depth, and flags as
ordinary search axes; parameter count and FLOPs reported as free
objectives. \textbf{One-shot supernets}: exclude (reliability evidence
against, RL mismatch). \textbf{Zero-cost proxy subsystem}: exclude
(free quantities capture most of the value). \textbf{Grammar-based
spaces}: exclude for now (solves expressiveness, not cost; no current
project needs topology discovery); NePS is the reference if that
changes.
\end{itemize}
\end{decisions}

\section{What would overturn these verdicts}

The routine-budget RL verdict (DEHB or BO+ASHA) is decided by V07;
whichever loses is excluded and this section rewritten. The sampler
study's veto gates are tied to V13: if vetoes prove unreliable
(failing configurations promoted, or survivors ranking poorly
against the reference posterior), correctness checking is pulled
back into training code and the tuner's veto hook is disabled. The
Hamiltonian study's reliance on the premium MO searcher follows
V09's outcome; if qLogNEHVI is demoted, the study defaults to
scalarization. The workload-template abstraction itself (four
matched studies, not a general tuner) holds unless colleagues
request a fifth workload family with evidence that it differs
structurally from the four, at which point a new template is argued
and added or the request is accommodated by parameterizing an
existing one.
